\bookStart{Introduction to Eddic Poetry}

The following two sections (Norse Mythic Poetry and Norse Heroic Poetry from the Walsing-Cycle) contain the bulk of preserved Old Norse Eddic poetry, most of which is preserved in only one manuscript, \Regius.  For this reason some notes about Eddic Poetry in general are in order.

\section{Character}

Old Norse poetry is conventionally divided into two categories: the Eddic and the Scaldic.

Eddic poetry is anonymous, composed in one of three meters (\Fornyrdislag, \Ljodahattr, \Malahattr), and free from regular assonance or end-rhyme.  It is easy to understand, uses kennings sparsely, and generally deals with heroic or mythological subjects.

By contrast, Scaldic poetry is usually attributed to named poets, is composed in a greater variety of meters (which apart from the three Eddic meters include the exclusively Scaldic \Kviduhattr\ and \Drottkvett), and employs regular assonance or end-rhyme.  It revels in purposeful obscurity by making great use of complex kennings and circumlocutions, and usually deals with historical or biographical subjects.  It typically serves as praise poetry for a living king or other chieftain.

Although this categorization is not explicitly made in any Norse primary source, it is not entirely a philological construct, for it is clearly observed in Snorre’s Edda.  In the first book, \Gylfaginning, we find almost exclusively citations made from Eddic poems (68 out of 70 sts.) while the second book, \Skaldskaparmal, has the inverse distribution (TODO).

\section{Dating}

Since Eddic poetry represents a mostly unified genre in a single language, it is useful to discuss some recent scholarly advancements in how this type of poetry can be dated.  Dating the original time of composition is important, since it allows us to tell whether the religious or cultural ideas in these texts do infact represent the views---indeed, the very words---of Germanic Heathens, or whether they are the recollections---in the worst case, the fantasies---of Christian antiquarians.

Before we discuss the dating, we must, however, clarify what we mean by a \emph{poem}.  We shall assume that the poems exist as authorial works independent of their manuscript form, that is, that a copy of a given poem in a certain manuscript represents a single written documentation of a poem, but is not synonymous or to be conflated with the poem itself.  The poem was composed as a single work by a single author; it can exist in multiple copies, and has most likely been transmitted both orally and scribally.  On the other hand, scribes make errors, and are not authors in their own right \parencites{Males2023Textual}{Males2024Emend}{Haukur2020}.  Even if they were, it is the primary goal of the present work to attempt get back to the wording of the original authors, be they Heathen or Christian, for it is due to \emph{them} alone that these texts fascinate us so much more than the contemporary output of the same scribes.

The ways of dating Eddic poetry can be divided into two categories.  External ways of dating rely on circumstances outside of the poems themselves.  These are citations in securely dateable texts, the lifetimes of the original authors or other documented circumstances of composition, and the dates of the individual manuscripts.  Unfortunately, the bulk of Eddic poetry is not cited in any texts older than Snorre’s Edda (composed ca.~1220); we have no surviving information about the original authors nor the compositional circumstances (beyond the attributed Greenlandish origin of \textlink{Atlakvida} and \textlink{Atlamal}); and even if we were able to use scribal characteristics to trace the origins of certain poems back to the very beginning of Icelandic vernacular scribal tradition---and Iceland is where all surviving Eddic manuscripts were produced---that tradition does not commence until the 12th century, which post-dates the time of conversion by over a century.

Internal ways of dating instead rely on the language, meter, allusions, and ideas of the poems themselves, and it is in these ways that the chief scholarly advancements of the last few decades have occurred.  These criteria can also allow us to distinguish different layers within a given poem.  For example, we can say that the Dwarf-tallies, \textlink{Voluspa}[11]--15, are later inserts into that poem.

To pinpoint the absolute date of when a certain linguistic or metrical change has occurred we require some external reference point, but, fortunately, this is found in the Scaldic corpus, which unlike the Eddic corpus is largely ascribed to named authors working within clear historical contexts.  It turns out that the Scaldic poems ascribed to the oldest poets likewise show the greatest concentration of archaic traits, and that some of these traits only occur in poems composed before certain dates.  Comparison with Scaldic poetry is therefore our best bet to date the Eddic corpus.
A landmark study in this direction is \textcite{Sapp2022}, which, although certainly not perfect, has generally been favourably received, and shown that fruitful results can be obtained at scale from applying linguistic dating criteria derived from Scaldic poetry to Eddic.

The most important linguistic criteria are the alliteration between etymological \emph{vr-} and \emph{v-} rather than \emph{r-} (pointing to composition before ca.~1000 \ce, \cite{Haukur2017}), the use of \emph{of} in prenominal positions analogous to West Germanic \emph{ga-} (pointing to composition before ca.~1030 \ce, \cite{Males2025Of}),
the frequency of the particle \emph{of} in prenominal or preverbal positions analogous to West Germanic \emph{ga-} (a trait strongly correlated with \emph{vr-}, \cite[54]{Haukur2016}),
the frequency of syntactical V2 violations in \Fornyrdislag\ (a trait correlated with other early features, \cite{Haukur2012}),
and the use of metrically required hiatus forms like \emph{sáa, fríi, séi} (pointing to composition before ca.~1150 \ce).  In addition to these, the forms of certain individual words (e.g. \emph{haufuð} for \emph{hǫfuð}, \emph{*Sig·vǫrðr} for \emph{Sig·urðr}, \emph{fjǫl} for \emph{fjǫlð}) can also point to an earlier date.

When it comes to content-based criteria, namely allusions to other texts and religious and cultural conceptions, these vary greatly from poem to poem and can hardly all be discussed here.  One notable criteria is, however, the use of pagan god-names.  It has been noted (TODO: cite) that the use of pagan kennings (e.g.~\Hakonarmal\ 8/3: \emph{Óðins veðr} ‘Weden’s storm \ken{battle}’) diminishes drastically in Scaldic poetry immediately following the conversion, and only rises in frequency with the Icelandic antiquarian revival of the 12th century.  This observation can, however, be qualified further: while the names of gods like \emph{Njǫrðr} ‘Nearth’ and \emph{Ullr} ‘Wolder’ remain in use in kennings after conversion, not a single dateable Scaldic stanza composed after 1000 \ce\ uses the name \emph{Óðinn} ‘Weden’ in a kenning of any type, nor indeed in any other way save explicitly condemning the god.  This is not the case for synonyms for Weden like \emph{Yggr} ‘Ug \name{= Weden}’, and so it seems that there was a unique concern, not to say fear, about using the name Weden in particular.

Of course, Scaldic poetry has different purposes than Eddic poetry, and it makes plenty of sense for an antiquarian poet to insert the god Weden into his own heroic Eddic composition.  Even so, it seems likely that little poetry containing this name would have been composed in the period intervening between the conversion of Iceland (1000 \ce) and the antiquarian revival (beginning in ca.~1150 \ce).  Therefore, a heroic poem with old linguistic features and this theonym must likely be quite old.
It is, of course, important that all content-based criteria be considered critically alongside linguistic traits.  We should not label poems late out of hand simply because they, for instance, portray pre-Christian gods in a humorous fashion (e.g.~\textlink{Harbardsljod}, \textlink{Lokasenna}), nor should the presence of an archaic archeological object like the ring-sword (\textlink{HelgakvidaHjorvardssonar}[9], \textlink{Sigurdskamma}[68]) make us conclude that the poem must go back to the Migration Period, when such swords were in use.  We cannot necessarily assume that pre-Christian poets never depicted the gods humorously (indeed, even \Iliad\ does), nor that a 12th century antiquarian poet could not have read about ring-swords even if he himself had not personally seen them.

I wish to conclude with two complementary examples of content-based and linguistic criteria working together.  When the term “burned” is used as a neutral synonym for “dead” in \textlink{Havamal}[71]/3, 81/1---a poem which also features prenominal \emph{of} and \emph{vr-}---it should be taken as a sign of pre-Christian composition, since cremation was a strong taboo in the post-conversion period \parencite[93--96]{Males2024Hav}.  Inversely, when \textlink{Atlamal}---a poem displaying the sound change \emph{vr-} > \emph{r-} and a very low occurrence of preverbal \emph{of}---links cremation to stoning as a severe punishment befitting an evil woman (a practice beloning to Judeo-Christian religion) and mentions Greenlandish fauna like polar bears, it clearly shows itself to be a late poem composed after the conversion to Christianity and the settlement of Greenland (both ca.~1000 \ce).

\section{Manuscripts}

\subsection{Codex Regius (R)}

By far the most important manuscript of Norse Eddic Poetry is GKS 2365 4to (siglum \Regius), the so-called Codex Regius.  It dates to around 1270 and consists of 45 surviving foll., containing 29 poems.  \Regius\ is fragmentary, however, and would originally have consisted of 53 foll.; see below.

\Regius\ ought properly to be divided into two parts or sections: the first, on foll. 1--20, contains 11 poems (\textlink{Voluspa}--\textlink{Grimnismal},  \textlink{Skirnismal}--\textlink{Alvissmal}) dealing chiefly with mythology; the second, on foll. 20--45, contains 18 poems (\textlink{HelgakvidaOne}--\textlink{Gudrunarhvot}) dealing with heroic legend from the Walsing cycle.
Distinct scribal and orthographic characteristics show that these two parts have been copied from two separate source manuscripts, likely by the scribe of \Regius\ itself \parencite[esp.~pp.~160--163]{Lindblad1980}.  Their distinct origin is supported by the characteristic that \Regius\ introduces the first poem in each part, \textlink{Voluspa} and \textlink{HelgakvidaOne}, with a particularly large initial \parencite[166]{Lindblad1980}.

If we are to look for relatives of these two source manuscripts, it seems likely that \AM, described below, is a collection of mythological poems of the same kind as that underlying the first half of \Regius.  Likewise, a manuscript closely related to the heroic half of \Regius\ has clearly served as the main source for large swathes of the younger \VolsungaSaga.

\Regius\ is not a mere anthology of poems, but shows substantial editorial input in the form of prose segments, likely accumulated over time.  A short prose section (\textlink{Lokasenna}[P1]) ties two of the mythological poems together into a loose overarching narrative, though it is clear from their meter, style, and language that they are are separate works composed by various poets over time.  When it comes to the heroic poems long prose segments occur both within and between them, creating a \inx[C]{saw}-like prosimetrical form where the prose sometimes comes to dominate the poetry.  The two long prose segments \textlink{GudrunTwo}[P1] and \textlink{Brot}[P1]--\textlink{GudrunOne}[P1] are often edited as separate works.

As noted above, a large gap famously occurs in the heroic half, the so-called Great Lacuna.  Between foll. 32 and 33 one quire, consisting of 8 foll., has gone missing \parencite[6 n. 2]{Lindblad1977}.  Its contents are mostly unknown, but it would at least have included the end of \textlink{Sigrdrifumal} and the beginning of \textlink{Brot}, alongside other now-lost poems or prosimetra linking together these two poems.  Some of the stanzas likely contained in the Lacuna may be restored from citations in \VolsungaSaga, and these are edited under \textlink{Lacuna}.

\subsection{AM 748 I a 4to (A)}

Second in importance stands AM 748 I a 4to (siglum \AM).  It dates to around 1300 and is in fragmentary state, consisting of just 6 foll.  The beginning and end are absent, and between foll. 2 and 3 there is a lacuna, so that at least 3 (but probably more) foll. are missing.

\AM\ contains seven poems.  On 1r--2v are found in succession the latter half of \textlink{Harbardsljod}, the full \textlink{Baldrsdraumar}, and the first half of \textlink{Skirnismal}.  There is then the lacuna---Finnur Jónsson guesses that just one fol. is missing---and on 3r--6v are found in succession most of \textlink{Vafthrudnismal}, all of \textlink{Grimnismal} and \textlink{Hymiskvida}, and the introductory prose to \textlink{Volundarkvida}.  Among mediæval mss., \textlink{Baldrsdraumar} is only attested in \AM, while the other six poems are also found in the first, mythological, half of \Regius. The order of the poems varies drastically between \AM\ and \Regius.

\AM\ has no trace of \Regius’s frame narrative tying together \textlink{Hymiskvida} and \textlink{Lokasenna} (and indeed the latter poem leaves no trace in \AM), but otherwise \AM\ and \Regius\ do share a substantial amount of prose (\textlink{Grimnismal}[P1]--P3, \textlink{Skirnismal}[P1], \textlink{Volundarkvida}[P1]), a fact which proves that the two, rather than being independent witnesses of poems based on oral tradition, stem from a common manuscript archetype.

\AM\ is written in the same hand and on the same parchment as \AMb\ (see below), a manuscript containing grammatical and poetic texts with which \AM\ was originally bundled.

\subsection{Manuscripts of Snorre’s Edda}

Snorre’s Edda or the Prose Edda consists of three sections.  The first two---\Gylfaginning\ and \Skaldskaparmal---contain quotations from several Eddic poems.  Snorre reproduces stanzas from the mythological \textlink{Voluspa}, \textlink{Vafthrudnismal}, \textlink{Grimnismal}, \textlink{Skirnismal}, \textlink{Lokasenna}, and \textlink{Hyndluljod} in \Gylfaginning\ and stanzas from \textlink{Alvissmal} in \Skaldskaparmal.  In addition he cites several otherwise unknown stanzas in Eddic meters, of which at least some appear to derive from now-lost mythological poems.  These fragments are all edited at the end of the present section under \textlink{EddicFragments}.

The four main mss. for the Prose Edda are:%TODO: use table like in the Heliand introduction

\begin{enumerate}
  \item \RegiusProse: Codex Regius of the Prose Edda (GKS 2367 4to), dating to 1300--1350 \ce.
  \item \Trajectinus: Codex Trajectinus (Traj 1374), a ca. 1595 paper copy of a ms. closely related to \RegiusProse.
  \item \Wormianus: Codex Wormianus (AM 242 fol.), dating to 1340--70. \Wormianus\ also contains the \textlink{Rigsthula} and the four Icelandic grammatical treatises.
  \item \Upsaliensis: Codex Upsaliensis (DG 11), dating to 1300--25.
\end{enumerate}

When all four mss. agree on a reading, the abbreviation \GylfMS\ is used synonymously with \RegiusProse\Trajectinus\Wormianus\Upsaliensis.

Of these mss.~\RegiusProse\ and \Trajectinus\ form their own stemmatic clade to the exclusion of \Wormianus\ and \Upsaliensis.  Stematically, \Upsaliensis\ is the most archaic ms., but unfortunately it has been heavily abbreviated and is filled with errors.  This makes it a questionable source, especially for poetry.  For detailed discussion on their internal stemmatics and origins I refer to \textcite{Haukur2017}.

Parts of \Skaldskaparmal\ are also preserved in three independent manuscripts:

\begin{itemize}
  \item \AMb: AM 748 I b 4° (\url{https://handrit.is/manuscript/view/is/AM04-0748-Ib})
  \item \EddaBms: AM 757 a 4° (\url{https://handrit.is/manuscript/view/is/AM04-0757a})
  \item \EddaCms: AM 748 II 4° (\url{https://handrit.is/manuscript/view/is/AM04-0748-II})
\end{itemize}

The first two of these also preserve the short poetic treatise known as the \emph{Lítla skálda} ‘Little Scalda’, which cites three mythological Eddic stanzas: \textlink{Grimnismal}[41]--42 and \textlink{EddicFragments}[F7][VIII] 1.

\AMb\ also contains the Third Grammatical Treatise and a copy of the \emph{Íslendingadrápa}, all written in the same hand.  \AMb\ is written in the same hand and on the same parchment as \AM, and it is clear that the two originally made up a single manuscript, something confirmed by the fact that they were found tied together into a codex alongside the unrelated \EddaCms.

TODO: \EddaBms\ and \EddaCms.

\subsection{Other manuscripts}

A few other Eddic-style poems from various sources are also included in the present section, both under Mythic Poetry.  As noted above, \textlink{Rigsthula} is attested in \Wormianus.  \textlink{Hyndluljod} is found in \FlatMS.  In addition, a younger redaction of \textlink{Voluspa} is found in \Hauksbok.

\subsection{Abbreviations}

The following abbreviations are relevant for the preservation of Eddic poems and will be used in the critical notes.  Some stanzas are cited on their own in unrelated texts, especially saws, preserved in various other manuscripts; these will be discussed in their individual places.

\begin{itemize}%
	\item \AM\ = AM 748 I a 4° (\url{https://handrit.is/manuscript/view/da/AM04-0748-I-a})
  \item \AMb\ = AM 748 I b 4° (\url{https://handrit.is/manuscript/view/is/AM04-0748-Ib})
  \item \EddaBms\ = AM 757 a 4° (\url{https://handrit.is/manuscript/view/is/AM04-0757a})
  \item \EddaCms\ = AM 748 II 4° (\url{https://handrit.is/manuscript/view/is/AM04-0748-II})
	\item \FlatMS\ = \emph{Flateyjarbók}, GKS 1005 fol. (\url{https://handrit.is/manuscript/view/is/GKS02-1005})
%	\item \GylfMS\ = all manuscripts of \Gylfaginning; equivalent to \RegiusProse\Trajectinus\Wormianus\Upsaliensis
	\item \Hauksbok\ = \emph{Hauksbók}, AM 544 4° (\url{https://handrit.is/manuscript/view/en/AM04-0544})
  \item \GylfMS\ = \RegiusProse\Trajectinus\Wormianus\Upsaliensis
	\item \VolsungaMS\ = NKS 1824 b 4°, the main ms.~of \VolsungaSaga\ (\url{https://onp.ku.dk/onp/onp.php?m9641})
	\item \Regius\ = Codex Regius of the Poetic Edda, GKS 2365 4° (\url{https://eae.ku.dk/q?p=eae/vols/text/1})
	\item \RegiusProse\ = Codex Regius of the Prose Edda, GKS 2367 4° (\url{https://handrit.is/manuscript/view/is/GKS04-2367})
	\item \Trajectinus\ = Codex Trajectinus, Traj 1374ˣ
	\item \Upsaliensis\ = Codex Upsaliensis, DG 11
	\item \Wormianus\ = Codex Wormianus, AM 242 fol. (\url{https://clarino.uib.no/menota/text/menota/AM-242-fol})
\end{itemize}
